\documentclass[11pt]{article}
\usepackage[margin=0.8in]{geometry}
\usepackage{amsmath,amssymb,mathpartir}
\usepackage[T1]{fontenc}
\usepackage{lmodern}
\usepackage[hidelinks]{hyperref}
\newcommand{\m}[1]{\operatorname{#1}}
\title{Operational Semantics}
\author{}
\date{}
\begin{document}
\maketitle
\vspace{-2em}

\section{Notation}
We retain $\Sigma=(M,P,C,S,K,T)$, with $K\in\mathbb{Q}_{>0}$ and
$T\in\mathbb{Q}_{\geq0}$. The ordinary source await is unstarted; its runtime
form $\text{await}^{\bullet}$ has an active timer. The existing frame grammar
includes $(tid:p,[];c)$ and $(tid:p,\text{let }x=[])$ for sequence and return
continuations. Expressions retain their small-step rules; $\rightarrow^*$
means complete expression evaluation, including zero steps for values.

$\m{bg}(S)(K)$ is the total background resource cost for one computational
step at clock $K$. At runtime $S$ contains all peripherals. Thus it includes
both the known background peripherals and those represented symbolically by
$\m{cost}(\cdot)$ at a function call site; the latter is not command overhead.
All resource labels below use the draft's clock-scaling convention.

\section{Expressions and sequential commands}
\begin{mathpar}
\inferrule*[left=E-Var]
{ M(x) = v }
{ M, S \vdash x \rightarrow v }

\inferrule*[left=E-BinopStep]
{ M, S \vdash e_1 \rightarrow e'_1 }
{ M, S \vdash e_1 \text{ binop } e_2 \rightarrow e'_1 \text{ binop } e_2 }

\inferrule*[left=E-BinopStep2]
{ M, S \vdash e_2 \rightarrow e'_2 }
{ M, S \vdash v_1 \text{ binop } e_2 \rightarrow v_1 \text{ binop } e'_2 }

\inferrule*[left=E-BinopFin]
{ v_1 \text{ binop } v_2 = v }
{ M, S \vdash v_1 \text{ binop } v_2 \rightarrow v }

\inferrule*[left=E-Seq]
{ q = 0 }
{
(M,P,(tid:p,c_1;c_2)::C,S,K) \rightsquigarrow_q
(M,P,(tid:p,c_1)::(tid:p,[];c_2)::C,S,K)
}

\inferrule*[left=E-SeqDone]
{ q = 0 }
{
(M,P,(tid:p,\text{skip})::(tid:p,[];c_2)::C,S,K) \rightsquigarrow_q
(M,P,(tid:p,c_2)::C,S,K)
}

\inferrule*[left=E-Call]
{
M,S \vdash e_a \rightarrow^* v \\
\m{body}(M(f),v,p) = (c';\text{ret}:=e) \\
C' = (tid:r,\text{let }x=[])::C \\ q = 0
}
{
(M,P,(tid:r,\text{let }x=f(e_a,p))::C,S,K) \rightsquigarrow_q
(M,P,(tid:r,c';\text{ret}:=e)::C',S,K)
}

\inferrule*[left=E-Ret]
{ M,S \vdash e \rightarrow^* v \\ C = (tid:p,\text{let }x=[])::C' \\ q = 0 }
{
(M,P,(tid:p,\text{ret}:=e)::C,S,K) \rightsquigarrow_q
(M[x\mapsto v],P,(tid:p,\text{skip})::C',S,K)
}

\inferrule*[left=E-IfTrue]
{
M,S \vdash e \rightarrow^* \text{true} \\
C' = (tid:p,c_1)::(tid:p,\texttt{endif})::C \\ q = \m{bg}(S)(K)
}
{
(M,P,(tid:p,\text{if }e\text{ then }c_1\text{ else }c_2)::C,S,K)
\rightsquigarrow_q (M,P,C',S,K)
}

\inferrule*[left=E-IfFalse]
{
M,S \vdash e \rightarrow^* \text{false} \\
C' = (tid:p,c_2)::(tid:p,\texttt{endif})::C \\ q = \m{bg}(S)(K)
}
{
(M,P,(tid:p,\text{if }e\text{ then }c_1\text{ else }c_2)::C,S,K)
\rightsquigarrow_q (M,P,C',S,K)
}

\inferrule*[left=E-EndIf]
{ q = 0 }
{
(M,P,(tid:p,\text{skip})::(tid:p,\texttt{endif})::C,S,K) \rightsquigarrow_q
(M,P,(tid:p,\text{skip})::C,S,K)
}

\inferrule*[left=E-Assign]
{ M,S \vdash e \rightarrow^* v \\ q = \m{bg}(S)(K) }
{
(M,P,(tid:p,\text{let }x=e)::C,S,K) \rightsquigarrow_q
(M[x\mapsto v],P,(tid:p,\text{skip})::C,S,K)
}
\end{mathpar}

The sequence rules use the continuation frame already present in the draft,
so calls and branches may push frames without losing the following command.
$\m{body}$ instantiates the stored function's formal parameters with $v,p$
capture-avoidingly and freshens its locals. Call entry and return copying are
free administrative steps, retaining the original convention. The expression
rules do not separately charge software instruction costs.

\section{Peripheral, clock, and asynchronous commands}
\begin{mathpar}
\inferrule*[left=E-Transition]
{ S(x) = st' \\ st' \leq st \\ q = \m{cost}(st',st)(K) + \m{bg}(S\setminus x)(K) }
{
(M,P,(tid:p,\text{transition}(x,st))::C,S,K) \rightsquigarrow_q
(M,P,(tid:p,\text{skip})::C,S[x\mapsto st],K)
}

\inferrule*[left=E-Downgrade]
{ S(x) = st' \\ (st'\hookrightarrow st)\in x(D) \\ q = \m{cost}(st',st)(K) + \m{bg}(S\setminus x)(K) }
{
(M,P,(tid:p,\text{downgrade}(x,st))::C,S,K) \rightsquigarrow_q
(M,P,(tid:p,\text{skip})::C,S[x\mapsto st],K)
}

\inferrule*[left=E-Clock]
{ K' = f \\ q = \m{cost}_{\text{clock}}(K,K') + \m{bg}(S)(K) }
{
(M,P,(tid:p,\text{clock }f)::C,S,K) \rightsquigarrow_q
(M,P,(tid:p,\text{skip})::C,S,K')
}

\inferrule*[left=E-Match-i]
{ S(p) = st_i \\ q = \m{bg}(S)(K) }
{
(M,P,(tid:r,\text{match }p\text{ with }st_1\rightarrow c_1|\dots|st_n\rightarrow c_n)::C,S,K)
\rightsquigarrow_q (M,P,(tid:r,c_i)::C,S,K)
}

\inferrule*[left=E-Async]
{
tid \notin \m{dom}(P)\cup\m{threads}(C) \\
\m{unstable}(S(p)) \\
q = C_{\text{spawn}}(p,S(p),K) + \m{bg}(S\setminus p)(K)
}
{
(M,P,(tid':r,\text{let }tid=\text{async}(p)\{c_a\})::C,S,K)
\rightsquigarrow_q
(M,P[tid\mapsto(p,c_a)],(tid':r,\text{skip})::C,S,K)
}

\inferrule*[left=E-AsyncPreempt]
{ P(tid) = (p,c_a) \\ tid \notin \m{threads}(C) \\ q = 0 }
{
(M,P,C,S,K) \rightsquigarrow_q (M,P,(tid:p,c_a)::C,S,K)
}

\inferrule*[left=E-AsyncDone]
{
C = (tid:p,\text{skip})::C' \\ tid \neq \text{main} \\
\m{root}(tid,C') \\ \m{stable}(S(p)) \\ q = 0
}
{
(M,P,C,S,K) \rightsquigarrow_q
(M,P[tid\mapsto\text{done}(p,\text{skip})],C',S,K)
}

\inferrule*[left=E-AwaitStart]
{ q = 0 }
{
(M,P,(tid':p,\text{await }tid\text{ timeout}[n]\{c_{\text{fb}}\})::C,S,K)
\rightsquigarrow_q
(M,P,(tid':p,\text{await}^{\bullet}\ tid\text{ timeout}[n]\{c_{\text{fb}}\})::C,S,K)
}

\inferrule*[left=E-Await]
{ P(tid) = \text{done}(p,\text{skip}) \\ q = 0 }
{
(M,P,(tid':r,\text{await}^{\bullet}\ tid\text{ timeout}[n]\{c_{\text{fb}}\})::C,S,K)
\rightsquigarrow_q
(M,P\setminus tid,(tid':r,\text{skip})::C,S,K)
}

\inferrule*[left=E-AwaitWait]
{ P(tid) = (p,c_a) \\ n>0 \\ q = \m{bg}(S)(K) }
{
(M,P,(tid':r,\text{await}^{\bullet}\ tid\text{ timeout}[n]\{c_{\text{fb}}\})::C,S,K)
\rightsquigarrow_q
(M,P,(tid':r,\text{await}^{\bullet}\ tid\text{ timeout}[n]\{c_{\text{fb}}\})::C,S,K)
}

\inferrule*[left=E-Await-Timeout]
{ P(tid) = (p,c_a) \\ tid\notin\m{threads}(C) \\ q = 0 \\ n\leq0 }
{
(M,P,(tid':r,\text{await}^{\bullet}\ tid\text{ timeout}[n]\{c_{\text{fb}}\})::C,S,K)
\rightsquigarrow_q
(M,P\setminus tid,(tid':r,c_{\text{fb}})::C,S,K)
}

\inferrule*[left=E-Timeout-Running]
{
P(tid) = (p,c_a) \\ n\leq0 \\
C = C_a\mathbin{+\!+}((tid':r,\text{await}^{\bullet}\ tid\text{ timeout}[n]\{c_{\text{fb}}\})::C_0) \\
\m{allFramesOf}(C_a,tid) \\ q = 0
}
{
(M,P,C,S,K) \rightsquigarrow_q
(M,P\setminus tid,(tid':r,c_{\text{fb}})::C_0,S,K)
}

\inferrule*[left=E-HWStep]
{ P(tid) = (p,\_) \\ st' \in \m{HWReach}(S(p)) \\ q = 0 }
{
(M,P,C,S,K) \rightsquigarrow_q (M,P,C,S[p\mapsto st'],K)
}
\end{mathpar}

\section{Global step and timer draining}
Let $r$ be the name of the local rule used in a derivation. Define
$\m{tick}(r)=1$ for
\textsc{E-Assign}, \textsc{E-IfTrue}, \textsc{E-IfFalse},
\textsc{E-Transition}, \textsc{E-Downgrade}, \textsc{E-Clock},
\textsc{E-Match-i}, \textsc{E-Async}, and \textsc{E-AwaitWait};
otherwise $\m{tick}(r)=0$. This is a classification of rules, not an extra
machine component or an inference from whether $q$ is zero.

\begin{mathpar}
\inferrule*[left=E-Step]
{
(M,P,C,S,K) \rightsquigarrow_q (M',P',C',S',K')\text{ by }r \\
C'' = \m{drain}_{\m{tick}(r)K}(C')
}
{
(M,P,C,S,K,T) \xrightarrow{q}
(M',P',C'',S',K',T+\m{tick}(r)K)
}
\end{mathpar}

Define $\m{drain}$ on frames and stacks:
\[
\begin{aligned}
D_d(tid:p,\text{await}^{\bullet}\ h\text{ timeout}[n]\{c\})
 &= (tid:p,\text{await}^{\bullet}\ h\text{ timeout}[\max(0,n-d)]\{c\}),\\
D_d(F)&=F\quad\text{for every other frame},\\
\m{drain}_d(\cdot)&=\cdot,\\
\m{drain}_d(F::C)&=D_d(F)::\m{drain}_d(C).
\end{aligned}
\]
Only active timers are decremented, including those below a running job.
Unexecuted awaits inside commands and continuations are unchanged.
Start a current await before dispatching async work. Process completion
bookkeeping and expired active timers before another computational step.
If completion has already been recorded, use \textsc{E-Await}.

\section{Declarations}
\begin{mathpar}
\inferrule*[left=E-PeriphDecl]
{ S' = S[id\mapsto s_{\text{curr}}] }
{
(M,S,id\{\langle s_0,\dots,s_n\rangle,s_{\text{curr}},<,D\})
\rightarrow (M,S')
}

\inferrule*[left=E-FunDecl]
{ M' = M[f\mapsto\{c;\text{ret}:=e\}] }
{
(M,S,\text{fn }f(x,p):(\tau_1,\vec s_1)\rightarrow(\tau_2,\vec s_2)\{c;\text{ret}:=e\})
\rightarrow(M',S)
}
\end{mathpar}
Peripheral metadata and formal parameter bindings remain available to the
rules; $S$ itself maps peripherals to their current states. Declarations are
elaborated before execution and are not charged to the program.

\section{Conventions needed by the rules}
\begin{itemize}
\item $\m{threads}(C)$ lists thread identifiers in the stack.
$\m{root}(tid,C')$ holds when $C'$ is empty or its top frame belongs to another
thread, preventing premature completion at an intermediate skip.
Handles are hygienically bound names. Local variable names are freshened;
async code cannot modify parent-local variables. Capture excludes access by
other software threads until completion is awaited or timeout hands off.
\item Transition costs include the changing peripheral's own resource use,
which explains $S\setminus x$. Likewise $C_{\text{spawn}}$ includes the
captured peripheral's own cost during spawning. If it does not, add that
peripheral's background cost to the spawn charge and its typing rule.
\item Hardware evolution is represented by instantaneous state changes with
resource use included in state/background costs, so \textsc{E-HWStep} does not
charge a second impulse cost. Exclude infinite zero-time hardware runs.
If separately charged hardware-transition costs are retained instead, their
sum needs a corresponding typing allowance.
\item Timeout retains the actual state $S$ and assumes an abstract exclusive
handoff to the fallback; its implementation is outside scope. The running
rule shown assumes that the only job frames above the waiting parent belong
to the awaited job (for example, one outstanding async job). Unrestricted
multi-job scheduling needs a more general handler rule. Check the fallback
from every possible state at expiry, not just the spawn state.
\end{itemize}

\paragraph{Timeout cost correspondence.}
Your expression
\[
B=C_{\text{unstable}}(p)+\m{cost}(\Delta\setminus p)+\m{cost}(\cdot)
\]
already covers partial activity if, interpreted as a rate, it satisfies
$q_{\text{step}}\leq KB$ for every computational step during active await.
There is no separate partial-work term. Explicit transition costs must also
be covered by this property when such commands are allowed.
For boundary-checked expiry with $K\leq K_{\max}$, the conservative bound is
$(n+K_{\max})B+Q_{\text{fb}}$; $nB+Q_{\text{fb}}$ requires exact deadline
handling or a normalization removing overshoot. Work before await activation
is accounted for outside that interval. Resource consumption of a run is
$Q_{\text{run}}=\sum_i q_i$; the global rule preserves these labels without
adding a resource counter to the machine.
\end{document}
